% Research in progress.
\subsection{A limit of the proposed four-point tests}

The following addresses the specific common-kernel route proposed in
\emph{UNSTUCK\_\allowbreak POINCARE\_\allowbreak FERMAT}, Section 4. It does not replace the
icosahedral membership problem by an assertion about its linearization.

\begin{statement}{C.23}{Four-point maps of the tetrahedral orbit family}{Proved}
Let $A_4\subset\operatorname{PGL}_2(\mathbb C)$ be the tetrahedral group,
let $S$ be its fourteen points with nontrivial stabilizer, and put
$T=\mathbb P^1(\mathbb C)-S$. The ordered orbit
$(g(z))_{g\in A_4}$ defines a map $T\to M_{0,12}$.
For every four-element subset $I\subset A_4$, let
$f_I:T\to M_{0,4}$ be the corresponding forgetful map. Then
\[
 \bigcap_{|I|=4}\ker\bigl(\widehat{\pi_1(T)}
                   \longrightarrow\widehat{\pi_1(M_{0,4})}\bigr)
 \ne1.
 \tag{C23.1}
\]
The kernel contains an explicitly specified nontrivial discrete word.
It therefore cannot be made trivial by passing from discrete groups to
their full profinite completions.
\end{statement}
\begin{proof}
There are six fixed points of involutions and two four-point orbits
fixed by subgroups of order three. These are disjoint, so $|S|=14$.
Choose peripheral loops $c_1,\ldots,c_{14}$ with
$c_1\cdots c_{14}=1$; the first thirteen form a free basis.

For four distinct fractional-linear functions $g_1,\ldots,g_4$, their
cross ratio is
\[
 \frac{(g_1-g_3)(g_2-g_4)}{(g_1-g_4)(g_2-g_3)}.
\]
After cancelling the four linear denominators, numerator and denominator
each have degree at most four. Thus a nonconstant $f_I$ has degree at
most four and at most twelve distinct preimages of $0,1,\infty$.
It maps $T$ to the tripod because equalities among orbit labels occur
only at $S$. At least two members of $S$ therefore map to interior
points of the tripod. The map extends across each such puncture and
kills its peripheral loop in fundamental groups. A constant $f_I$
kills every loop.

Use the convention $[a,b]=aba^{-1}b^{-1}$ and set
\[
 w_2=[c_1,c_2],\qquad w_j=[w_{j-1},c_j]
 \quad(3\le j\le14).
 \tag{C23.2}
\]
Filling any one of the fourteen punctures kills $w_{14}$. Hence every
$f_I$ kills it. Its nontriviality is not inferred from this formula
alone: the exact free reduction in
\texttt{root\_crossratio\_kernel.py}, after substituting
$c_{14}=(c_1\cdots c_{13})^{-1}$, leaves a word of length $24\,594$.
The attached finite permutation certificate has degree $24\,595$ and
sends $0$ to $24\,594$ under this word. To construct that certificate,
follow the reduced word along a path with distinct successive vertices,
assign each signed edge to the corresponding generator permutation,
and complete each resulting partial injection to a permutation.
Reducedness excludes conflicting assignments. This also proves
nontriviality in the full profinite completion directly.

The program independently enumerates the group generated by
$z\mapsto-z$, $z\mapsto1/z$, and $z\mapsto(z+i)/(z-i)$.
Its exceptional affine polynomial is
\[
 z^{13}+13z^9-13z^5-z,
\]
with infinity as the fourteenth point. Among all $\binom{12}{4}=495$
cross ratios, $351$ have degree four and $144$ have degree two.
Every one kills at least two boundary loops, as the degree argument
requires. The exact coefficient calculations and the two permutation
certificates are retained with their output.
\end{proof}

The discrete orbit monodromy is injective. Indeed the quotient by $A_4$
is a sphere with three branch points and one moving puncture; the
ordered orbit family is its twelve-sheeted parameter cover. A loop
with trivial upstairs mapping class projects to the identity downstairs
by \citet[Theorem 1.1]{Winarski2015}: the regular branched cover has no
unramified preimage of a branch point, and the upstairs twelve-punctured
sphere has negative Euler characteristic. Downstairs the moving-point
group is $\pi_1(M_{0,4})$, so the loop is trivial. Thus the word above
is also a nontrivial pure mapping class. Pure sphere mapping classes
embed by their action in $\operatorname{Out}(F_{11})$ and act trivially
on its abelianization; \citet[Theorem 1.4]{Paris2009} separates any
nonidentity such element in a finite group. This verifies that the
specific discrete witness remains nontrivial in the ambient profinite
mapping-class group. We do not infer that the entire completion of the
parameter group embeds merely from its discrete injectivity.

\begin{statement}{C.24}{The cyclic-five family has the same defect}{Proved}
Fix a primitive fifth root $\zeta$. On
$T_5=\mathbb P^1-\{0,\infty,1,\zeta,\zeta^2,\zeta^3,\zeta^4\}$,
the ordered configuration
\[
 0,\ \infty,\ 1,\ \zeta,\ldots,\zeta^4,
 \ z,\ z\zeta,\ldots,z\zeta^4
\]
has twelve distinct points. The common kernel of all its four-point
forgetful homomorphisms on $\widehat{\pi_1(T_5)}$ is nontrivial.
The corresponding nontrivial discrete monodromy also survives in the
ambient profinite mapping-class group.
\end{statement}
\begin{proof}
Every label is constant or linear in $z$, so each pairwise projective
determinant has degree at most one. Every cross ratio has degree at
most two, and hence uses at most six of the seven removed points as
preimages of $0,1,\infty$. It kills at least one peripheral loop.
The exact enumeration in \texttt{root\_c5\_crossratios.py} finds
$70$ constant maps, $325$ maps of degree one, and $100$ of degree two;
it checks all $495$ subsets over $\mathbb Q(\zeta)$.
With $c_7=(c_1\cdots c_6)^{-1}$, use the nested commutator (C23.2)
through $j=7$. Its reduced length is $196$; the attached permutation
certificate of degree $197$ sends $0$ to $196$. Filling any peripheral
kills the word, so it belongs to every four-point kernel.

For injectivity of the discrete monodromy, forget all variable labels
except $z$. The seven constant labels remain, and the parameter curve
is exactly the fibre of $M_{0,8}\to M_{0,7}$. The homotopy exact
sequence gives injection on fundamental groups: the base has zero
$\pi_2$, by iterating the forgetful fibrations down to $M_{0,3}$,
whose fibres are punctured spheres and are aspherical. The preceding
residual-finiteness argument applies to its pure mapping-class image.
\end{proof}

Thus the proposed common-kernel test cannot prove icosahedral rigidity:
its kernel is already nontrivial in each of these discrete symmetry
families. Nor do these words produce an element of $\mathcal K$.
Being killed by all four-point maps is weaker than lying in the actual
intersection of the two family subgroups inside the full mapping-class
group. In particular, a tetrahedral word and the identity in the cyclic
family pass every such cross-ratio comparison while representing
different full mapping classes. This locates the precise information
lost by that proposed reduction, and leaves the actual membership tests
of C.20 as the remaining problem.
